%======================================================================
\section{Introduction}
\label{sec:intro}
%======================================================================

Zero-shot forecasters are trained once and applied to datasets they have never seen. The open models
that do this well, among them Moirai \citep{moirai}, Chronos \citep{chronos}, Time-MoE
\citep{timemoe} and TimesFM \citep{timesfm}, reach that ability by pretraining on time-series corpora of tens to
hundreds of billions of observations. The cost is the point: it decides who can build such a model
and how often it can be rebuilt when the data distribution moves. We ask what the accuracy ceiling
looks like at a small fraction of that budget, and we get there by starting from a backbone that
already knows something about how signals evolve.

The visual route shows the idea is viable. VisionTS \citep{visionts} renders a series as an image and
lets a frozen ImageNet masked autoencoder \citep{mae} inpaint the missing portion; VisionTS++
\citep{visiontspp} continues that backbone's pretraining on rendered series with a multi-quantile
head and reports the strongest published zero-shot numbers on the standard long-sequence protocol.
Both inherit a prior over spatial structure from a model that has never seen anything change over
time. Video pretraining offers a different inheritance: representations built from sequences and, in the
predictive variants, an explicit forward model \citep{videomae, vjepa2}. The useful question about
it is not whether it raises the ceiling but whether it lowers the cost of reaching one: \emph{does a video-pretrained backbone reach a given forecasting accuracy under a smaller
time-series adaptation budget?}

Getting an answer takes more than swapping the backbone. A numeric series has to become a video whose
axes carry temporal meaning; masked reconstruction, the objective the backbone arrives with, recovers
a token from neighbours on both sides, which is interpolation rather than forecasting; and the
backbone minimises a gray level while the task scores a number, with an untrained decoder between
them. The consequence is measurable: continuing VideoMAE's native objective on a 654M-observation
corpus of rendered series lowers its held-out pixel loss monotonically and raises forecast error
monotonically across the checkpoints we evaluated, $0.339$ at 5k updates to $0.381$ at 176k
(Section~\ref{sec:diagnosis}). More of that pretraining makes the forecaster worse.

We therefore keep the video weights and rebuild the adaptation stage around forecasting
(Figure~\ref{fig:overview}). \textbf{Period-aligned rendering} puts one period on the horizontal
axis of a frame and value on the vertical axis, so that the frame axis is time measured in periods
and a video backbone's extrapolation axis lines up with the series' own. \textbf{Future-block
prediction with a differentiable read-out} masks a trailing block of frames, predicts it from the
visible past only, decodes the predicted frames back into values, and takes the loss on those
values, which puts the pixel-to-number decoder inside the optimisation instead of after it.
\textbf{Temporal-scale alignment} trains on exactly the periods-per-frame scales the inference rule
selects, so the model is never asked at test time for a resolution it was not trained at. These are
design choices for a transfer problem, not three independent inventions: Section~\ref{sec:related}
places each against its nearest existing method, and Section~\ref{sec:design} reports the ablation
that each is supported by.

The result is a position on the trade-off between accuracy and resources rather than a large
accuracy jump.
Sixty thousand optimizer updates at an effective batch of 32 over 654M observations, which is 18.9
hours on one A40, about $0.3\%$ of the LOTSA archive and $1/27$ of the sample presentations behind
VisionTS++, give one checkpoint that, under one inference rule, reaches a mean MSE of $0.285$ over six
standard datasets and four horizons: the lowest of the models we compare (VisionTS++ large $0.289$,
base $0.291$, VisionTS $0.309$), by a small margin. On mean MAE we are behind, $0.328$ against
$0.326$, and per cell against VisionTS++ base we take 17 of 24 on MSE but 11 on MAE. The efficiency
statement is the substantive one; the accuracy statement is that we stay in the leading band while
making it.

Two controlled experiments locate the video prior's contribution. Changing only the starting weights,
it leads an ImageNet initialisation by $3$ to $5\%$ at every horizon at 5k updates, and by 20k the
image initialisation has caught up; separately, removing the pretrained weights altogether costs
$14\%$ under an otherwise identical trainer. What video pretraining supplies here is adaptation
efficiency, the binding constraint in this regime, and not a higher ceiling. We ran that
comparison at 5k and 20k updates and not at the 60k of the reported system, so the efficiency claim
is the one the study supports; we do not claim the headline number depends on the video
initialisation.

\paragraph{Contributions.}
\begin{itemize}
\item A transfer framework that makes a video-pretrained backbone usable as a zero-shot forecaster:
      period-aligned video rendering, future-block prediction supervised through a differentiable
      numeric read-out, and a training scale distribution matched to the inference rule
      (Section~\ref{sec:method}).
\item Evidence that this reaches the accuracy band of the strongest published zero-shot models at
      $0.3\%$ of their time-series corpus, $1/27$ of their sample presentations and 18.9 GPU-hours,
      with the budget reported in separated units and no compute claim
      (Sections~\ref{sec:main} and~\ref{sec:budget}).
\item A measurement of what the video prior is worth and where it stops: a matched-size, matched-recipe
      initialisation study with a random-init control at two adaptation budgets
      (Section~\ref{sec:backbone}).
\item A diagnosis of why the naive route fails, the backbone's own objective and forecasting
      accuracy moving in opposite directions as pretraining continues, which motivates the framework
      and cautions anyone continuing a vision objective on rendered series
      (Section~\ref{sec:diagnosis}).
\end{itemize}
